\documentclass[11pt]{article}

\usepackage{amsmath}
\usepackage{amssymb}
\usepackage[margin=1in]{geometry}
\usepackage{booktabs}
\usepackage{graphicx}
\usepackage{hyperref}

\newcommand{\zref}{\hat{z}_{\rm ref}}
\newcommand{\khat}{\hat{k}}
\newcommand{\lhat}{\hat{l}}
\newcommand{\rhat}{\hat{r}}
\newcommand{\nhat}{\hat{n}}

\title{Polarization conventions in the Monte Carlo module}
\author{Athena++ Monte Carlo radiation transport}
\date{}

\begin{document}
\maketitle

\begin{abstract}
This note fixes the conventions the Monte Carlo module uses to report polarization: which
frame the output is measured in, which axis the Stokes parameters are referenced to, what
the sign of $Q$ means, and how the two independent angle-binning choices relate to them.
It is a companion to \texttt{doc/monte\_carlo/frames\_and\_moments.tex}, which covers the
frames themselves and the radiation moments. The emphasis here is on the conventions rather
than the transport, because almost every question about a polarized result turns out to be
a question about which plane $Q$ was measured against.
\end{abstract}

\section{The output frame}

Polarization is reported in the frame of the normal (Eulerian) observer at the photon's
position,
\begin{equation}
  n^\mu = -\alpha \, g^{\mu t} , \qquad \alpha = \frac{1}{\sqrt{-g^{tt}}} ,
\end{equation}
which stays well defined inside an ergosphere. \texttt{NormalFrameWavevector} builds it and
projects the photon's wavevector into it; \texttt{ConstructTetrad} supplies the tetrad.

The orientation of the three spatial legs is worth stating explicitly, because it is
inherited rather than chosen. \texttt{ConstructTetrad} takes the coordinate directions
$\delta^\mu_{\ 1}, \delta^\mu_{\ 2}, \delta^\mu_{\ 3}$ as trial vectors and orthonormalises
them by Gram--Schmidt \emph{in index order}, each projected orthogonal to $e_{(0)}$ and to
the legs already built. The legs are therefore tied to the coordinate ordering of the chart:

\begin{table}[h]
\centering
\begin{tabular}{llll}
\toprule
chart & $e_{(1)}$ & $e_{(2)}$ & $e_{(3)}$ \\
\midrule
cartesian $(x,y,z)$          & $\hat{x}$   & $\hat{y}$      & $\hat{z}$ \\
spherical polar $(r,\theta,\phi)$ & $\hat{e}_r$ & $\hat{e}_\theta$ & $\hat{e}_\phi$ \\
cylindrical $(R,\phi,z)$     & $\hat{e}_R$ & $\hat{e}_\phi$ & $\hat{e}_z$ \\
\bottomrule
\end{tabular}
\end{table}

The third column is the one that matters. In a cartesian or cylindrical chart the third leg
\emph{is} the global $z$ axis, so nothing distinguishes ``the tetrad's third leg'' from
``the global $z$ axis''. In spherical polar the third leg is $\hat{e}_\phi$, which rotates
with azimuth.

Because the frame is orthonormal and the photon is null, the projected components satisfy
$|k_{\rm spatial}| = k^{(0)}$. Output code relies on that.

\section{The meridian plane}

Given the unit propagation direction $\nhat = \khat$ in that frame and a reference axis
$\zref$, \texttt{MeridianPair} constructs
\begin{equation}
  \lhat = \frac{\zref - (\zref \cdot \nhat)\,\nhat}{\left| \zref - (\zref \cdot \nhat)\,\nhat \right|} ,
  \qquad
  \rhat = \nhat \times \lhat ,
  \label{eq:meridian}
\end{equation}
so that $(\lhat, \rhat, \nhat)$ is a right-handed orthonormal triad. The \emph{meridian
plane} is the plane spanned by $\zref$ and $\khat$. By construction $\lhat$ lies in it and is
perpendicular to the propagation direction, while $\rhat$ is normal to it
(Fig.~\ref{fig:meridian}).

\begin{figure}[h]
\centering
\includegraphics[width=0.78\textwidth]{meridian_plane.pdf}
\caption{The meridian plane and the basis Stokes parameters are referenced to. The plane is
spanned by the reference axis $\zref$ and the propagation direction $\khat$. Removing the
component of $\zref$ along $\khat$ and normalising gives $\lhat$, which lies in the plane and
is perpendicular to $\khat$; $\rhat = \nhat \times \lhat$ is normal to the plane. $Q>0$ means
the polarization lies along $\lhat$, in the plane; $Q<0$ means it lies along $\rhat$,
perpendicular to it. Generated by \texttt{meridian\_figure.py}, which asserts the
orthonormality and handedness it draws.}
\label{fig:meridian}
\end{figure}

If $\nhat$ is parallel to $\zref$ the meridian is undefined. \texttt{MeridianPair} returns
false, and the caller returns \emph{without updating}: the Stokes parameters silently retain
whatever values they already held rather than being zeroed. This is worth remembering when
reading photons that emerge close to the reference axis.

\section{Which axis is the reference}

The reference axis in Eq.~\eqref{eq:meridian} is $(0,0,1)$ in tetrad components --- the
tetrad's third spatial leg --- except in spherical polar, where the global $z$ direction is
projected into the tetrad instead, using
\begin{equation}
  \frac{\partial}{\partial z} = \cos\theta \, \frac{\partial}{\partial r}
                              - \frac{\sin\theta}{r} \, \frac{\partial}{\partial \theta} .
\end{equation}

The reason is the table in Sec.~1. Left on the third leg, $Q$ and $U$ in spherical polar
would be referenced to a meridian containing $\hat{e}_\phi$, a direction that rotates with
azimuth. Every photon's $Q$ would then be measured against a different physical plane, the
values could not be combined across the grid, and --- the sharpest symptom --- the emergent
polarization would not vanish for a photon leaving along the polar axis, where by symmetry it
must.

With that projection in place, all four combinations agree on the global $z$ axis:

\begin{table}[h]
\centering
\begin{tabular}{lll}
\toprule
pusher & chart & reference axis \\
\midrule
legacy  & cartesian & global $\hat{z}$ (the transport convention, see \texttt{polarization.hpp}) \\
legacy  & spherical & global $\hat{z}$ (same) \\
general & cartesian & global $\hat{z}$ (the third leg already is) \\
general & spherical & global $\hat{z}$ (projected in, as above) \\
\bottomrule
\end{tabular}
\end{table}

Only the last row required anything; before it was added, that case was the odd one out.

\section{Stokes parameters}

\texttt{WriteMeridianStokes} projects the coherency tensor into the tetrad, contracts its
spatial block onto the meridian pair to give $N_{ll}, N_{rr}, N_{lr}, N_{rl}$, and hands
those to \texttt{TensorToStokes}. Following Mo\'scibrodzka \& Gammie (2018),
\begin{equation}
  I = \tfrac{1}{2}\left( N_{ll} + N_{rr} \right) , \qquad
  Q = \tfrac{1}{2}\left( N_{ll} - N_{rr} \right) , \qquad
  U = \tfrac{1}{2}\left( N_{lr} + N_{rl} \right) , \qquad
  V = \tfrac{1}{2}\,\mathrm{Im}\left( N_{rl} - N_{lr} \right) .
\end{equation}

So \textbf{$Q>0$ means the polarization lies along $\lhat$}, in the meridian plane, and
$Q<0$ means it lies along $\rhat$, perpendicular to it. $U$ is the component at $\pm 45$
degrees between $\lhat$ and $\rhat$.

Two consequences catch people out.

\paragraph{They are stored normalised.} \texttt{TensorToStokes} divides through by $I$, so
the stored $I$ is identically one and \texttt{sqp}, \texttt{sup}, \texttt{svp} are the
\emph{fractions} $Q/I$, $U/I$, $V/I$. The absolute intensity enters the outputs separately,
as $w \times E$. This is why \texttt{PhotonList::AddPhoton} writes only $Q$, $U$ and $V$ and
never $I$.

\paragraph{The physically expected sign is negative.} A scattering atmosphere polarizes
perpendicular to the plane containing the surface normal and the line of sight --- which is
exactly the meridian plane when $\zref = \hat{z}$. Chandrasekhar's classical result for a
semi-infinite electron-scattering atmosphere, polarization vanishing at $\mu=1$ and rising to
$11.7\%$ at $\mu = 0$, therefore appears here as $Q<0$. The Feautrier reference in
\texttt{vis/python/montecarlo/feautrier.py} reports $Q/I$ in this same convention, so the
comparison scripts use the Monte Carlo output directly with no sign change.

\section{Where the referencing happens}

\texttt{CoherencyToObserverStokes} refreshes the stored Stokes parameters from the
transported coherency tensor. It runs in the terminal sweep of
\texttt{TransferPhotonsOnBlock}, \emph{before} \texttt{FinalizePhoton}, before
\texttt{Spectrum::UpdateSpectrum} and before \texttt{PhotonList::AddPhoton}. Every output
therefore sees the referenced values.

It returns immediately for the legacy pushers, which maintain $Q$ and $U$ against the global
cartesian meridian throughout the transport and need no refreshing: in flat spacetime a
photon travels a straight line, $\khat$ is constant in cartesian components, so the meridian
plane is constant and the Stokes parameters do not change between scatterings.

Note that the referencing above is for \emph{output}. The scattering routines use a different
convention --- the comoving frame and that frame's own third leg --- because
\texttt{ScatterThomsonPolarized} reads the reference off the stored components directly. The
two are deliberately separate, and \texttt{MeridianBasis} passes $(0,0,1)$ explicitly for
that reason.

\section{Angle binning}

Two independent choices control the angle $\mu$ a spectrum is binned on, and they have
opposite defaults.

\begin{table}[h]
\centering
\begin{tabular}{lll}
\toprule
 & option & default \\
\midrule
internal \texttt{spec} output & \texttt{<output$n$>/polar\_axis} & \texttt{false}: bins on $\khat\cdot\hat{e}_r$ \\
\texttt{make\_spectrum.py}    & \texttt{--anglebin}              & \texttt{cartesian}: bins on $|k_z|$ \\
\bottomrule
\end{tabular}
\end{table}

For a disk, the meaningful angle is against the global $z$ axis: $\mu = 1$ looks down the
rotation axis and $\mu \to 0$ grazes it edge-on. That is \texttt{polar\_axis = true} for the
internal output and the default for \texttt{make\_spectrum.py}. Binning on the global axis is
what makes the reference axis of Sec.~3 the right one to have chosen: the two must agree, or
the $\mu$ dependence of $Q$ is meaningless.

Two details. \texttt{make\_spectrum} bins on $|k_z|$, the \emph{absolute value}, so the two
hemispheres are folded together; that is harmless when only one surface radiates but will
co-add them when both do. And \texttt{make\_spectrum} accumulates the stored $Q$ and $U$
verbatim --- it does not reconstruct the meridian --- so it inherits whatever referencing the
run used.

\section{A known inconsistency}

In spherical polar, the photon list writes $k$ in \emph{tetrad} components while $Q$ and $U$
are referenced to the global $z$ axis. These are different bases, which contradicts the
comment in \texttt{PhotonList::AddPhoton} that the wavevector ``goes out in the same normal
frame the Stokes parameters are referenced to'', and the statement in
\texttt{NormalFrameWavevector} that a reader can rebuild the meridian basis from the stored
components alone by taking the third spatial leg as $(0,0,1)$.

Nothing in the repository currently depends on that guarantee --- \texttt{make\_spectrum}
uses $Q$ and $U$ verbatim and rotates $k$ into cartesian itself --- but any consumer written
against the documented invariant would get the wrong basis. Closing it means either rotating
the list's $k$ into global cartesian for spherical as well, so the two frames agree again, or
amending both comments to state that the axes differ. The first is the more useful invariant;
it also changes the $k$ columns of existing spherical photon lists, so it is not free.

\end{document}
